\section{Method}
\label{sec:method}

\subsection{Overview: two connected execution paths}
\label{sec:overview}

The system contains a shared deterministic data substrate and two deliberately distinct
language-model interfaces.

The \textbf{typed-program path} maps a question directly to either an algebra tree or an
abstention object.  It is used to generate program supervision and to measure compositional
semantic-parsing capability.  A tree is validated, executed over the query-record universe, and
scored against an offline oracle.  This path does not create an evidence bundle or call the
deployment repair loop.

The \textbf{full runtime} first elicits a structured intent.  A router then chooses a
deterministic compiler whenever the intent can be compiled without either of two implemented
fast-path vetoes; remaining questions are sent to a typed graph planner.  Both branches converge
on a graph/low-level specification that is grounded, preflighted, executed, and checked before
an answer is released.  A bounded feedback path is available only for eligible structured
defects.  Figure \ref{fig:method-overview} is reserved for this general control flow, whereas
Figure \ref{fig:visual-thesis} must contain one concrete trace.

\begin{figure}[t]
  \centering
  \figureplaceholder{1.65in}{\textbf{Author-drawn Figure 2: full runtime.}\\
  Question $\rightarrow$ structured understander $\rightarrow$ routing gate
  $\rightarrow$ deterministic intent compiler \emph{or} typed graph planner
  $\rightarrow$ consistency/compile $\rightarrow$ grounding and preflight
  $\rightarrow$ deterministic execution $\rightarrow$ evidence, postflight, and sanity checks
  $\rightarrow$ answer with evidence / no verified answer.  One dashed bounded-replan edge.}
  \caption{General full-runtime control flow.  A directly compilable intent bypasses the graph
  model; a structured failure may trigger bounded replanning, after which the complete
  ground--execute--check sequence is repeated.  The separate typed-program evaluation omits
  these deployment stages.}
  \label{fig:method-overview}
\end{figure}

\subsection{From OCDS releases to a queryable graph}
\label{sec:data-method}

The ingestion stage scans yearly gzip-compressed JSON Lines files.  It skips empty and malformed
lines rather than terminating a multi-year build, unwraps release packages, and flattens the
fields needed by downstream extraction.  OCDS publishes a new release when information about a
contracting process changes \citep{ocds-standard}; consequently, multiple rows can share an Open
Contracting ID (OCID).  For each OCID, the implementation retains the row with the latest
parseable release date and records all observed source years.  This policy creates a current
snapshot rather than a complete event history.  It is deterministic for a fixed input ordering
and explicitly differs from merging all release versions field by field.

The full pipeline entry point accepts an optional year subset.  In the current implementation,
partial runs must be directed to an explicit non-canonical output path; this prevents a small
diagnostic run from silently replacing a complete artifact.  The canonical full build preserves
its original default path.

Entity resolution is conservative because a false merge affects every downstream join.  Phase 1
uses official organisation identifiers when an observed identifier belongs to an accepted
official scheme.  A publication-system identifier (such as a GB-FTS identifier) can become an
alias only when observations with the same OCID and normalised name share a party role and the
complete candidate set points to exactly one official canonical identifier.  Candidate
identifiers are accumulated across observations before the uniqueness decision; a first-row
match is not sufficient.  This rule avoids merging two same-named parties that occur as buyer
and supplier or that map to different official organisations.

Phase 2 resolves remaining parties in three deterministic passes: a government lookup table;
exact normalised name plus region; and, when region is absent, exact normalised name.  A resolved
cluster receives a stable identifier derived from SHA-256.  An observation whose name
normalises to an empty string is retained as its own singleton instead of disappearing from the
output.  Jaro--Winkler similarity at a threshold of 0.92 is used only to emit review candidates;
it never performs an automatic fuzzy merge.  Thus the resolver favours false negatives over
irreversible false positives.

The extraction stage produces normalised tables for tenders, lots, award criteria, awards, bid
statistics, text evidence, and documents.  For an award-level monetary value, the implementation
uses the award value when present, otherwise a linked contract value, and finally a tender value
as a fallback.  Only award- and contract-sourced values are labelled
\texttt{value\_is\_additive}; tender fallbacks are not.  A numeric zero is a valid value and is
not replaced by a later fallback.  Currency is stored but not normalised.

Every extracted fact retains identifiers that link it to the OCID, release, award or contract,
and source evidence where available.  These links permit the runtime to return record-level
provenance, but they do not prove that a publisher supplied every real-world fact.

The graph contains organisation, contract/notice, CPV, and evidence nodes.  Buyer and supplier
edges connect canonical organisations to contract nodes; category edges connect contracts to
CPV nodes; and evidence edges connect source artifacts to contracts.  Contract-node identifiers
are deterministically formed from the OCID and award identifier.  Before publication, the build
checks primary-key uniqueness, referential integrity, JSON serialisability, additivity labels,
and coverage thresholds.  The frozen experimental graph contains \KGContracts contract nodes,
\KGOrgs organisation nodes, \KGCPV CPV nodes, and \KGEvidence evidence nodes.  Exact edge and
coverage counts appear in Table \ref{tab:dataset-stats}.

The query backend materialises a flat record universe $U_G$ from these Parquet nodes and edges.
This representation intentionally makes the unit of \texttt{count} explicit: duplicate rows are
deduplicated by \texttt{contract\_node\_id}.  Buyer/supplier fields use the corresponding stored
name projection.  Grouping drops null and empty keys.  This record view is the contract between
the algebra and the graph; changing it changes benchmark semantics and therefore requires a new
snapshot hash.

\subsection{Closed typed algebra and deterministic execution}
\label{sec:algebra}

The algebra is \emph{closed} in its operator inventory but permits open composition of those
operators.  Static validation limits a tree to depth 16 and 64 nodes and returns one of seven
root types.  Table \ref{tab:operators} summarises the principal signatures; Appendix
\ref{app:algebra} provides the complete inventory and failure semantics.

\begin{table}[t]
  \caption{Core algebra operators.  $P$ is a predicate, $f$ a field, $m$ an aggregate metric,
  and $\bowtie_s$ an expression-valued membership predicate.}
  \label{tab:operators}
  \centering
  \small
  \begin{tabular}{lll}
    \toprule
    Family & Signature & Semantics \\
    \midrule
    Filter & $\operatorname{filter}_{P}:\records\!\to\!\records$ & Select records \\
    Project & $\operatorname{values}_{f}:\records\!\to\!\values$ & Distinct non-empty values \\
    Cardinality & $\operatorname{count}:\records\!\to\!\numbertype$ & Distinct contract count \\
                & $\operatorname{size}:\values\!\to\!\numbertype$ & Set cardinality \\
    Aggregate & $\operatorname{sum}_{f}:\records\!\to\!\numbertype$ & Decimal accumulation \\
    Group & $\operatorname{groupby}_{f,m}:\records\!\to\!\groups$ & Keyed count or sum \\
    Set & $\operatorname{union/intersect/diff}:\values^2\!\to\!\values$ & Set algebra \\
    Bridge & $\operatorname{filter}_{f\bowtie_s e}$ & Semi-/anti-join with $e:\values$ \\
    Rank & $\operatorname{argext}:\groups\!\to\!\valuetype$ & Extreme group key \\
         & $\operatorname{top}_k:\groups\!\to\!\ranking$ & Ordered key--value pairs \\
    Scalar & $\operatorname{combine}:\numbertype^2\!\to\!\numbertype/\booltype$ & Arithmetic/comparison \\
    Boolean & $\operatorname{exists}:\records\!\to\!\booltype$ & Non-empty test \\
    \bottomrule
  \end{tabular}
\end{table}

Predicates support equality, membership, containment, existence, bounded numeric comparison,
negation, disjunction, and expression-valued membership.  The last form is the key to a bridge
query.  For example, if $k$ is an anchor buyer name,
\begin{align}
 S &= \operatorname{values}_{\texttt{supplier\_name}}
      (\operatorname{filter}_{\texttt{buyer\_name}=k}(U_G)),\\
 B &= \operatorname{values}_{\texttt{buyer\_name}}
      (\operatorname{filter}_{\texttt{supplier\_name}\in S}(U_G)),\\
 A &= B\setminus
      \operatorname{values}_{\texttt{buyer\_name}}
      (\operatorname{filter}_{\texttt{buyer\_name}=k}(U_G)).
 \label{eq:bridge-example}
\end{align}
The types make two common errors visible before execution: filtering
\texttt{supplier\_name} by a buyer anchor is a role error detectable only from semantics, but
returning contract identifiers where a set of buyer values is expected also violates the
declared answer signature.

Static validation checks tree shape and type composition, not the truth of the program or every
domain guard.  Field existence is checked during evaluation.  Monetary training templates add
an explicit \texttt{value\_is\_additive=true} predicate, but the generic \texttt{sum} operator
does not itself prove that this predicate is present, and neither evaluator currently enforces a
single currency.  We therefore treat monetary results as subject to the limitation in Section
\ref{sec:limitations}.

\label{sec:execution}
The runtime algebra evaluator first calls the static validator, then evaluates the tree
recursively over $U_G$.  A record filter constructs a Boolean mask.  An
\texttt{in\_expr} predicate first evaluates its child to a distinct value set and then applies
membership to the parent field.  Counts deduplicate contract identifiers.  Distinct projection
removes null and empty strings.  Numeric sums deduplicate records and accumulate with
\texttt{Decimal}; non-finite or non-numeric values fail the execution.  Grouped counts
deduplicate $(\text{group key},\text{contract id})$ pairs.  Rankings use a deterministic
value-and-key ordering, and select operations fail rather than guess when more than one distinct
answer remains.

A second evaluator is implemented separately for dataset generation and audit.  A candidate
training instance is retained only when the system evaluator and this independent implementation
both return successful envelopes and their answers agree under the generator comparator.  This
reduces circularity between template generation and execution, although it is not formal
verification: the two implementations can share assumptions, and the historical comparator had
an overly permissive Boolean coercion rule (Section \ref{sec:limitations}).

\subsection{Program-first supervision}
\label{sec:program-supervision}

Training examples are sampled from the actual record distribution rather than generated as
unanchored text.  A sampler selects concrete buyers, suppliers, years, CPV codes, categories, and
other field values; instantiates a typed program; renders a natural-language instruction; and
executes the program.  We reject a sample if (i) the tree fails static validation, (ii) either
evaluator fails, (iii) the evaluators disagree, (iv) the answer is degenerate for its declared
type, or (v) the question string has already been accepted.  Degeneracy rules exclude zero-valued
numeric examples, empty value sets, singleton rankings, and missing scalar answers.  These rules
are a data-construction heuristic, not a statement that zero is intrinsically invalid at
inference time.

The resulting LLaMA-Factory export contains \TrainN training and \ValN validation examples.
Each target is JSON containing either a typed tree or, for unsupported evaluation controls, an
abstention object.  The specialised parser adapts Qwen3-8B \citep{yang2025qwen3} with 4-bit QLoRA
\citep{dettmers2023qlora}; the optimiser settings are given in Section
\ref{sec:implementation}.  The held-out synthetic loss is reported only as an optimisation
diagnostic because the export contains duplicates and two exact train/validation overlaps.

\subsection{Understanding, conditional planning, and grounding}
\label{sec:understanding-planning}

The full runtime begins by retrieving a compact schema context and requesting a structured intent
program.  The intent records the question shape, answer signature, source record constraints,
intermediate entity-set roles, target constraints, and return operation.  Provider-enforced JSON
schema output is used when available, with a text/JSON fallback for older clients.  A malformed
briefing receives one cheap structural retry; an unparseable result becomes a typed planning
failure rather than a factual answer.

If the object is a recognised intent program, a deterministic compiler maps it to a graph plan.
The compiled fast path is accepted unless compilation fails or one of two narrow semantic vetoes
fires: an unsupported cue that could be converted into an invented answer, or a mismatch between
an entity-set count and a record-count question.  The implementation intentionally does not run
the raw-planner consistency checker over every compiled graph, because checks calibrated for raw
model payloads produced false positives on deterministic compiler output.

Questions that do not pass this gate invoke the Step-2 graph planner.  The planner emits typed
variables, dependencies, operation units, and a return specification.  A consistency checker
tests the payload against the question, and the compiler normalises and attaches low-level query
specifications.  Structural compile/consistency failures may trigger a small number of fresh
plan samples.  An executable but semantically wrong plan is not detectable from structure alone
and is not resampled using an oracle at inference time.

This conditional design is related to high-precision grammar plus neural fallback systems
\citep{shaw-etal-2021-compositional}.  A complete routing ablation would compare compiler-only,
always-planner, and conditional modes under the same model and budget.  The present artifact set
does not contain that complete three-way experiment, so the paper describes the implemented
mechanism but does not attribute the entire end-to-end gain to routing.

\label{sec:grounding}
Grounding maps aliases to canonical backend fields, normalises values, resolves confident
organisation mentions, and rejects fields outside the allowed backend profile.  It also attaches
operation-specific guards such as the additivity predicate used for monetary operations.  A
preflight phase checks schema support and query preconditions before any answer is produced.

For a graph plan, variables are executed in dependency order.  A source record-set variable
applies its grounded constraints.  An entity-set variable emits distinct values from its source.
The emitted set is then bound to the appropriate field of a dependent target variable---for
example, supplier values are bound to \texttt{supplier\_name}, whereas a record-set
intersection binds contract identifiers.  Variables at the same dependency level may execute in
parallel; a terminal operation computes the requested count, set, comparison, or other supported
answer.  Each variable trace stores its status, grounded constraints, emitted field, output size,
and answer summary.  The executor is the sole factual authority: the planner is never asked to
invent the terminal value.

\subsection{Release checks and bounded repair}
\label{sec:verification}

Execution returns record and evidence identifiers in an \texttt{EvidenceBundle}.  Postflight
checks inspect operation-specific properties of the returned rows, while answer-sanity checks
detect incompatible shapes, non-finite numbers, and other release defects.  An evidence verdict
summarises whether the implemented support requirement is satisfied.  A released
\texttt{AnswerCard} contains the denotation or its verbalisation, confidence label, evidence
identifiers, and limitations.  If a required check fails, the runtime produces no verified
answer.

These checks provide operational guarantees, not an oracle.  They can establish that a program
uses allowed fields, executes over the declared scope, and returns linked records.  They cannot
establish that the plan captures the user's intended semantics or that the publisher's data are
complete.  We therefore reserve \emph{correct} for the offline oracle metric and use
\emph{runtime verified} only in the sense defined by Table \ref{tab:guarantee-boundary}.

\label{sec:repair}
Algorithm \ref{alg:runtime} summarises the implemented control logic.  Repair is eligible when a
structured defect supplies actionable feedback: compilation or consistency rejection, failed
preflight/postflight checks, an execution/verifier failure, or a no-result plan containing a
literal not traceable to the question.  The feedback object records the failure stage, failed
checks, grounding information, and invented literals.  A new plan is then grounded, executed,
and checked from the beginning.  No repaired answer is accepted merely because its surface form
looks better.

Conversely, if a fully grounded query executes to no results and every literal is traceable to
the question, the empty result is terminal.  Replanning in this case would amount to
answer-shopping.  Unsupported operations, unresolved ambiguity, exhausted repair budgets, and
insufficient evidence all terminate with a controlled non-answer.

\begin{figure}[t]
  \centering
  \fbox{\begin{minipage}{0.93\linewidth}
  \small
  \textbf{Algorithm 1: Full question-answering runtime}\par
  \begin{enumerate}
    \item Obtain structured intent $i\leftarrow\textsc{Understand}(q,\mathrm{schema})$.
    \item If $i$ compiles and no fast-path veto fires, set $H\leftarrow\textsc{Compile}(i)$;
          otherwise obtain and consistency-check $H\leftarrow\textsc{Plan}(q,i)$.
    \item For each attempt up to the fixed budget: ground all variables and literals; run
          preflight; execute variables in dependency order; construct evidence; run postflight,
          sanity, and evidence checks.
    \item If all release conditions pass, return an answer card whose factual value is the
          executor denotation.
    \item If the defect is repair-eligible, send structured feedback to the planner and return
          to Step 3.  Otherwise return the corresponding unsupported, ambiguous, no-result, or
          insufficient-evidence non-answer.
  \end{enumerate}
  \end{minipage}}
  \caption{Runtime control procedure.  Repair changes the plan, never the factual answer
  directly, and repeats the complete deterministic path.}
  \label{alg:runtime}
\end{figure}
